% SOURCE_AUTHORITY: sga5_fr_workpass.tex, lines 13906--13951 (LF-normalized unit).
% SOURCE_SHA256: 791F4EFFC5E02832D5D77ED03518C8156D6F07E4C8238B03545DB93D883FBB28
\subsection*{6.3. Uma fórmula de Lefschetz}
Sob as hipóteses e com as notações de 5.1, 6.1, 6.2, tem-se a fórmula:
\begin{equation}
\tag{6.3.1}
\Tr_A(\RG_X(u,f))=\sum_{x\in X^f}\tau_x(f,u,f',G,\varphi),
\end{equation}
em que se pôs
\begin{equation}
\tag{6.3.2}
\tau_x(f,u,f',G,\varphi)=\Tr_A(u_x)-\alpha_x(f,u,f',G,\varphi).
\end{equation}

\begin{prooftext}[Demonstração]
Como o subesquema $Y$ é estável sob $f$, o morfismo
\[
(u,f):(X,F)\longrightarrow(X,F),
\]
induz um endomorfismo de triângulos distinguidos:
\begin{equation}
\tag{6.3.3}
\begin{tikzpicture}[baseline=(current bounding box.center),>=stealth]
\node (Y) at (0,1.75) {$\RG_Y(F|Y)$};
\node (U) at (-2.45,0) {$\RG_X(F_{U,X})$};
\node (X) at (2.35,0) {$\RG_X(F)$};
\node at (2.75,2.25) {$\scriptstyle\RG_Y(u|Y,f|Y)$};
\node at (-5.05,0.02) {$\scriptstyle\RG_X((u,f)_{U,X})$};
\node at (5.05,0.02) {$\scriptstyle\RG_X(u,f)$};
\draw[->] (Y) to[out=145,in=35,looseness=5.8] (Y);
\draw[->] (U) to[out=210,in=-30,looseness=5.5] (U);
\draw[->] (X) to[out=25,in=-25,looseness=7] (X);
\draw[->] (Y) -- (U);
\draw[->] (U) -- (X);
\draw[->] (X) -- (Y);
\end{tikzpicture}
\end{equation}
daí resulta a relação:
\begin{equation}
\tag{6.3.4}
\Tr_A(\RG_X(u,f))=
\Tr_A(\RG_X((u,f)_{U,X}))+
\Tr_A(\RG_Y(u|Y,f|Y)).
\end{equation}
O último termo da fórmula (6.3.4) pode ser explicitado mediante a fórmula de Lefschetz elementar e o fato de que podemos sempre supor $X^f=Y^f$:
\begin{equation}
\tag{6.3.5}
\Tr_A(\RG_Y(u|_Y,f|_Y))=\sum_{x\in X^f}\Tr_A(u_x).
\end{equation}
